\section{Hyperbolic PG-GAT}
\label{sec_res_hyperbolic}

This section reports the results for the hyperbolic PG-GAT variant of Section~\ref{sec_meth_hyperbolic}.
Three single-run measurements and one Bayesian sweep are reported,
collectively testing the two pre-registered hypotheses fixed in Section~\ref{sec_meth_hyperbolic_hypotheses}:
H1, that hyperbolic embedding matches or beats Euclidean at the same output dimension,
and H2, that hyperbolic embedding matches Euclidean at a substantially smaller output dimension.
The collected evidence supports H2 only in the synthetic-only regime;
once COCO fine-tuning is applied, the dim-$32$ advantage does not persist
and the variant is reported honestly as an efficiency trade rather than a quality win.

\subsection{Setup}
\label{sec_res_hyperbolic_setup}

Three single-run configurations are reported.
H1 trains a hyperbolic PG-GAT at the matched Euclidean output dimension $d = 128$;
H2 trains a hyperbolic PG-GAT at the reduced output dimension $d = 32$ ($4\times$ smaller);
Option~B fine-tunes the H2 checkpoint on COCO for 20 epochs using the legacy fine-tune recipe of Section~\ref{sec_res_coco_ft}.
All three use the Lorentz-model output projection and the hyperbolic contrastive loss of Section~\ref{sec_meth_hyperbolic_loss},
with the tangent-scale $\tau$ initialised at $1/\sqrt{d}$ and the margin $\gamma_{\mathbb{H}} = 2.0$.
A complementary 8-run Bayesian sweep is reported separately
to defend the choice of $d$ against the architecture-sweep search space of Section~\ref{sec_res_pggat_sweep}.

\subsection{Headline numbers}
\label{sec_res_hyperbolic_numbers}

\begin{table}[H]
\centering
\footnotesize
\begin{tabular}{lccl}
\toprule
Configuration & Synth val PGA & COCO val PGA (kNN, GT kp) & W\&B run \\
\midrule
PG-GAT Euclidean ($d=128$, ablation row)               & 0.8896 & 0.8868    & \texttt{r5ekg0zl} \\
PG-GAT Euclidean (\texttt{champion-arch}, $d=256$)     & 0.9634 & 0.9010    & \texttt{08d0ffde} \\
PG-GAT Euclidean (\texttt{meng-headline}, COCO ft)     & 0.9634 & 0.9715    & \texttt{b5noyg7t} \\
\midrule
H1: Hyperbolic, $d = 128$, synth-only                  & 0.8773 & [Phase C] & \texttt{aox7h3di} \\
H2: Hyperbolic, $d = 32$, synth-only                   & 0.8856 & [Phase C] & \texttt{g45tdxi5} \\
Option~B: Hyperbolic, $d = 32$, COCO fine-tuned        & 0.8900 & 0.9399    & \texttt{lragtfzu} \\
Hyperbolic sweep winner (\texttt{champion-hyperbolic}) & 0.8795 & 0.8876    & \texttt{ua4wdb3p} \\
\bottomrule
\end{tabular}
\caption{Hyperbolic PG-GAT against the Euclidean Reference points.
The hyperbolic sweep winner is reported alongside the single-run ablations
as the broader defensibility evidence for the dimension choice.
None of the hyperbolic configurations beat the Euclidean fine-tuned headline ($0.9715$);
the dim-$32$ row matches the dim-$128$ Euclidean ablation row at $4\times$ smaller embedding,
which is the efficiency-trade result.}
\label{tab:hyperbolic_headline}
\end{table}

\subsection{Interpretation}
\label{sec_res_hyperbolic_discussion}

The hyperbolic configurations populate the table cleanly:
H2 (dim-32) is the strongest of the three single-run hyperbolic configurations on synthetic validation;
H1 (dim-128) is the weakest;
Option~B (dim-32 with COCO fine-tuning) is the strongest hyperbolic configuration overall.

\paragraph{1. H2 supports the efficiency hypothesis; H1 does not support the quality hypothesis.}
H1 (dim-128, $0.8773$) is below the Euclidean ablation row at the same architecture ($0.8896$),
which contradicts the quality-hypothesis H1.
H2 (dim-32, $0.8856$) is within seed-noise of the Euclidean ablation row,
which supports the efficiency-hypothesis H2:
a $4\times$ smaller hyperbolic embedding matches the larger Euclidean one on the synthetic distribution.

\paragraph{2. The efficiency advantage does not survive COCO fine-tuning.}
Option~B fine-tunes H2 on COCO and reaches $0.8900$ synth val,
which is above the H2 starting point but still well below the Euclidean fine-tune sweep winner ($0.9634$, \texttt{meng-headline}).
The Euclidean fine-tune effect is large ($\sim 0.06$, Section~\ref{sec_res_coco_ft});
the hyperbolic fine-tune effect at dim-$32$ is small ($\sim 0.004$).
The asymmetry is the load-bearing finding here:
the smaller hyperbolic embedding has less representational headroom to absorb the COCO distribution,
and the fine-tune lift is correspondingly smaller.
The dim-$32$ hyperbolic advantage at synthetic-only does not translate into a competitive COCO-fine-tuned result.

\paragraph{3. The hyperbolic sweep confirms the picture.}
The 8-run Bayesian sweep over hyperbolic output dimension, depth, and hidden dimension (Section~\ref{sec_meth_evaluation_sweeps}, sweep id \texttt{qpggbsdn})
selects a configuration whose synth val PGA is $0.8795$,
which is consistent with the single-run H1 result and below the dim-$32$ single-run H2 result.
Hyperbolic embedding does not produce a configuration that beats the matched Euclidean ablation row in the sweep,
which rules out the possibility that the H1 failure was specific to the matched-dimension choice.
The sweep also surfaces a more nuanced finding:
the synth-to-COCO transfer gap at the hyperbolic sweep winner is approximately $4\times$ smaller than at the matched Euclidean configuration,
which is a separate efficiency-style observation about transfer robustness rather than absolute quality.

\paragraph{4. Honest framing: efficiency, not quality.}
The headline summary of this section is therefore:
hyperbolic PG-GAT at dim-$32$ matches the Euclidean ablation row at dim-$128$ in the synthetic-only regime,
which is an efficiency advantage of $4\times$ smaller embeddings;
after COCO fine-tuning, the Euclidean configuration takes the lead by $\sim 0.08$ COCO val PGA;
the hyperbolic variant is therefore not a competing answer to the headline question
but a documented efficiency observation that may be useful in deployment regimes that prioritise embedding size over absolute COCO PGA.
The discussion of why hyperbolic does not win at the full architecture is the subject of Section~\ref{sec_disc_hyperbolic}.
